\section{Fixed-trace reflection and proof-producing finite programs}
\label{sec:trace}

\subsection{State machines specialized at the theorem}

Fix a query tuple \(q\), a state code \(s\), and a transition function
\[
  T(s,q).
\]
All are represented by fixed primitive-recursive symbols.  Define the trace by primitive recursion:
\[
  \mathsf{run}(0,q)=s_0(q),
  \qquad
  \mathsf{run}(t+1,q)=T(\mathsf{run}(t,q),q).
\]
The current library's \code{recDef} gives the symbol and object-theory zero/successor equations at arbitrary terms.

A state may pack several logical components: loop variables, accumulators, a failure bit, table cursors, certificate digits, and resource counters.  Pairing or mixed-radix coding turns the tuple into one natural number.  Fixed-dimensional tuples do not require arrays as primitive objects.

\subsection{Quantifier-free invariants}

Let \(I(s,q)\) be a quantifier-free formula or, preferably, an invariant bit equation \(\iota(s,q)=1\).  Assume object derivations
\[
\begin{aligned}
&\Gamma(q)\vdash_{\PRAAx} I(s_0(q),q),\\
&I(s,q),\Gamma(q)^{+}\vdash_{\PRAAx} I(T(s,q),q).
\end{aligned}
\]
The fresh variable in the step sequent is the induction counter/state parameter as arranged by \code{IPRA.ind}.

\begin{theorem}[Fixed-trace reflection]
\label{thm:trace-reflection}
For every fuel term \(N(q)\), iPRA derives
\[
  \Gamma(q)\vdash_{\PRAAx} I(\mathsf{run}(N(q),q),q).
\]
If an object derivation extracts a terminal formula \(F(q)\) from that invariant, then \(\Gamma(q)\vdash_{\PRAAx}F(q)\).
\end{theorem}
\begin{proof}
Form the open induction formula
\[
  \phi(t,q):=I(\mathsf{run}(t,q),q).
\]
It is quantifier-free because \(I\) is and \(\mathsf{run}\) is a term.  The base case reduces by the zero equation of \(\mathsf{run}\) to the supplied base derivation.  The successor case reduces by the successor equation to the supplied step derivation.  Apply \code{IPRA.ind}, then instantiate the induction variable with \(N(q)\) by \code{Prf.subst}.  Compose the terminal extractor by simultaneous cut or implication elimination.
\end{proof}

\subsection{Booleanizing prefix properties}

An informal loop invariant often says ``all entries processed so far satisfy \(P\).''  The literal formula \(\forall i<t\,P(i)\) is not quantifier-free.  SCOPE instead extends the state with
\[
  a(0)=1,
  \qquad
  a(t+1)=a(t)\wedge_b \bit_P(t).
\]
The invariant contains the equation \(a(t)=1\).  After induction, bounded reflection identifies it with the intended prefix property.  A failure accumulator is dual:
\[
  f(0)=0,
  \qquad
  f(t+1)=f(t)\vee_b\neg_b\bit_P(t).
\]

This pattern handles table verification, finite covers, path checks, proof traces, and bounded universal side conditions.  It also makes counterexamples extractable by least search.

\subsection{Adaptive finite interactions}

A consumer may ask its next finite question based on previous answers.  Such adaptivity is still a fixed trace if:

\begin{enumerate}
\item the interaction horizon is bounded by a primitive-recursive term;
\item every answer lies in an explicit finite code range;
\item the next-query function is primitive recursive;
\item validity of each answer is reflected by a bit.
\end{enumerate}

The state stores the transcript.  The transition decodes the current answer, verifies it, updates the transcript and next query, and accumulates failure.  A private master certificate can encode the entire answer table in advance; the trace merely reads the demanded row.  This ``offline certificate, online verifier'' view is often easier to formalize than a genuinely interactive object.

\subsection{Phase composition}

Large algorithms are best split into phases with separate invariants.  Suppose phase \(j\) proves a postcondition \(P_j\), and phase \(j+1\) has a base derivation from \(P_j\).  Compile each phase by \cref{thm:trace-reflection} and compose the resulting derivations with \code{Prf.trans}.  There is no need for one enormous invariant carrying all earlier implementation details.

Phases also localize certificate scopes.  A certificate used only to establish \(P_j\) can be existentially eliminated before phase \(j+1\), preventing irrelevant data from polluting the public theorem and subsequent state.

\subsection{Termination and descent}

A loop whose natural description is ``repeat until done'' must expose primitive-recursive fuel.  Typical sources are:

\begin{itemize}
\item a finite search bound;
\item a decreasing natural rank;
\item a lexicographic tuple encoded by a rank;
\item a finite state-space cardinality together with a no-repeat certificate;
\item the capped demand budget of \cref{sec:demand-closure}.
\end{itemize}

The repository already contains generic descent developments.  SCOPE uses them when the step relation is theorem-specific, but retains the same discipline: termination evidence is a finite term and its correctness is an object derivation.

\subsection{A restricted proof-producing DSL}

The fixed-trace theorem suggests a small implementation language with nodes for:

\begin{itemize}
\item exact term computation;
\item reflected conditionals;
\item bounded folds and least search;
\item clocked recursion with an invariant;
\item certificate read/verify operations;
\item local existential scopes;
\item residual producer calls.
\end{itemize}

A compiler for this DSL need not be trusted for logical correctness.  It emits \(\PR\) symbols, formulas, and a \(\Prf\,\PRAAx\) term assembled from checked lemmas.  Lean checks the output.  The DSL is deliberately weaker than arbitrary Lean, which is what makes proof production feasible.
